\section{Related Work}


\paragraph{Single-Agent Approaches.} Recent advances in large language models (LLMs) such as GPT-4 \cite{openai2023gpt4} have renewed interest in the development of autonomous agents that can solve tasks on behalf of people \cite{russellnorvig1996, Wooldridge1995,Jennings1998Applications,wu2023autogen,yang2023autogptonlinedecisionmaking,sodhi2024stepstackedllmpolicies,zhang2024cognitivekernelopensourceagent,reed2022generalist}. These modern agents have shown remarkable skills in software development \cite{wang2024opendevinopenplatformai,zhang2024autocoderover,yang2024swe,xia2024agentless}, web manipulation \cite{deng2023mind2webgeneralistagentweb,zhang2024webpilotversatileautonomousmultiagent,zhou2024webarenarealisticwebenvironment,nakano2021webgpt,abuelsaad2024agenteautonomouswebnavigation}, manipulation of general graphical user interfaces \cite{zhang2024ufouifocusedagentwindows,wu2024oscopilotgeneralistcomputeragents,bonatti2024windowsagentarenaevaluating,pan2024autonomousevaluationrefinementdigital}, and other domains \cite{park2023generativeagentsinteractivesimulacra, Wang2023survey}. 
%Moreover, a subset of such agentic systems exhibit some generalist ability across different domains, however, they are not at a consistent state-of-the-art level across all domains \cite{wang2024opendevinopenplatformai,zhang2024cognitivekernelopensourceagent}.


Common strategies for developing such agents \cite{liu2024agent,xi2023risepotentiallargelanguage, masterman2024landscape, cheng2024exploring} include equipping LLMs with tools such as for code execution and web browsing \cite{qin2023tool,qin2023toolllm,schick-arxiv2023,mialon2023gaiabenchmarkgeneralai} and prompting strategies for better reasoning and planning such as CoT \cite{wei2022chain}, ReACT \cite{yao-iclr2023} and few-shot prompting \cite{zhou2024webarenarealisticwebenvironment}. With the development of multimodal models, agents can also operate in visual domains with techniques such as Set-of-Marks prompting \cite{yang2023set} among others \cite{yang2023set, zhang2024lookscreensmultimodalchainofaction,paranjape2023art,he2024webvoyagerbuildingendtoendweb}. To allow agents to accomplish tasks that require multiple steps with improved reliability, agent systems can incorporate self-critique \cite{wu2024oscopilotgeneralistcomputeragents,pan2024autonomousevaluationrefinementdigital,paul-etal-2024-refiner}, and inference-time search \cite{chen2024treesearchusefulllm,yao2023treethoughtsdeliberateproblem,koh2024treesearchlanguagemodel,song2024trialerrorexplorationbasedtrajectory}. Finally, Agentic systems can also benefit from memory and training either through explicit fine-tuning \cite{zeng2023agenttuningenablinggeneralizedagent,pan2024autonomousevaluationrefinementdigital,liu-arxiv2024,putta2024agentqadvancedreasoning} or through memory mechanisms \cite{wang2024agentworkflowmemory,sodhi2024stepstackedllmpolicies}.  Our work incorporates a subset of these techniques, and distributes them across agents in \team's multi-agent workflow, resulting in a modular, easy-to-extend implementation. 


\paragraph{Multi-Agent Approaches.} The multi-agent paradigm presents an attractive modular and flexible approach to tackling complex tasks \cite{MASaSurvey2000, Messing2002AnIT, Grosz1999SharedPlans,Tambe1998AgentTeams, Scerri2001AdjustableAI, wu2023autogen, talebirad2023multiagentcollaborationharnessingpower,guo2024large,liu2024agent,xi2023risepotentiallargelanguage,masterman2024landscape,cheng2024exploring}. Commonly each agent either has access to different tools or has a different role in the team, sometimes defined through the system prompt of the LLM or by explicit training. Sibyl presents a multi-agent approach with a debate-based jury mechanism with tools for python code execution and web browsing \cite{wang2024sibylsimpleeffectiveagent}. WebPilot uses a multi-agent system with global and local optimization  in planning for web based tasks \cite{zhang2024webpilotversatileautonomousmultiagent}. Trase claims to use a multi-agent architecture with a top level agent with self-critique and lower level agents \cite{redcell_trase_2024}.  A host of other multi-agent systems and frameworks have also been introduced \cite{li2023camel,liang-arxiv2023,du2023improving,babyagi,hong2023metagpt}. However, the previous methods differ from the architecture of \team which incorporates dynamic routing between agents using the Orchestrator along with planning and recovery. 




\paragraph{Agentic Evaluation.} To evaluate agents on general multi-step tasks, numerous benchmarks have been proposed in the literature \cite{mialon-arxiv2023,zhou2024webarenarealisticwebenvironment,xie2024osworldbenchmarkingmultimodalagents,liu2024visualwebbenchfarmultimodalllms,yoran2024assistantbenchwebagentssolve,yao2023webshopscalablerealworldweb,shi2017world,deng2023mind2webgeneralistagentweb,pan2024webcanvasbenchmarkingwebagents,li2024websuite}. Given the general and ubiquitous nature of the web, many of these benchmarks heavily incorporate \cite{mialon2023gaiabenchmarkgeneralai, yoran2024assistantbenchwebagentssolve}, or exclusively consider \cite{zhou2024webarenarealisticwebenvironment, deng2023mind2webgeneralistagentweb} browser-based tasks. These benchmarks either rely on non-interactive traces through real websites such as Mind2Web \cite{deng2023mind2webgeneralistagentweb}, interaction with synthetically created websites such as in WebArena \cite{zhou2024webarenarealisticwebenvironment}, or interaction with real websites on the public Internet such as GAIA \cite{mialon2023gaiabenchmarkgeneralai}. In the former case, non-interactive benchmarks are  limiting for evaluating agentic systems since they do not allow agents to deviate from previously recorded paths. This makes it impossible to evaluate error recovery, or find novel alternative strategies for the given problem.
Therefore, we focus on benchmarks that rely on interacting with live websites -- whether synthetic or public -- as they are more faithful to real-world tasks. Moreover, we prioritize benchmarks such as GAIA, which test generalist skills like data analysis or coding, in addition to commanding web browsers to navigate pages. We contribute \agbench{} as a standalone tool to perform evaluation of agentic systems, relying on benchmarks from the literature. Furthermore, we provide an in-depth error analysis of \team's performance contributing to work on debugging agentic systems \cite{li2024websuite}.

